% Introduction to inclusive breakup theory

Inclusive breakup reactions, where a projectile $a$ breaks into a detected fragment $b$ and an undetected particle $x$ that may fuse with the target or escape, sum over all final states of the $x+A$ system. The inclusive cross section $d^2\sigma/d\Omega_b dE_b$ is experimentally accessible even when the residual system is not fully resolved. It is the main observable for experiments with weakly bound projectiles at stable and radioactive beam facilities and underpins the surrogate method, where $(d,p)$ reactions proxy for neutron capture on short-lived nuclei.
\
\
The theoretical description of inclusive breakup was developed along three distinct lines in the early 1980s. Austern and Vincent [Phys.\ Rev.\ C 23, 1847 (1981)] wrote a post-form DWBA sum rule that expresses the inclusive cross section as an elastic-distorted-wave expectation of an optical-model propagator for the unobserved $x+A$ system. Ichimura, Austern, and Vincent [Phys.\ Rev.\ C 32, 431 (1985), IAV] extended this work and demonstrated the formal equivalence of the post and prior sum rules under a consistent treatment of nonorthogonality between entrance and breakup channels. Udagawa and Tamura [Phys.\ Rev.\ C 24, 1348 (1981), UT] independently derived a prior-form breakup-fusion formula from the optical theorem, stressing its link to the imaginary part of the optical potential, and applied it to inclusive spectra across many systems. Hussein and McVoy [Nucl.\ Phys.\ A 445, 124 (1985), HM] adopted the spectator-core model with closure to obtain a $\langle \psi | W | \psi \rangle$ estimate for the nonelastic breakup cross section, yielding a transparent eikonal-limit expression that became the standard tool for knockout reactions at intermediate energies. The coexistence of the three formulations sparked an extended debate centered on the post-prior ambiguity and on whether different formalisms counted the same physical mechanisms or included unphysical contributions [Ichimura et al., Phys.\ Rev.\ C 34, 2326 (1986); Udagawa et al., Phys.\ Rev.\ C 37, 2261 (1988); Ichimura et al., Phys.\ Rev.\ C 37, 2264 (1988)].
\
\
The post-prior equivalence was settled by Lei and Moro [Phys.\ Rev.\ C 92, 061602 (2015)], who performed the first comprehensive numerical test of the IAV post and prior sum rules for deuteron and $^6$Li-induced inclusive breakup reactions and confirmed their equivalence at the cross-section level. They extended the proof to transfer reactions with complex binding potentials [Lei and Moro, Phys.\ Rev.\ C 97, 011601 (2018)], showing that a complex final-state potential introduces an additional term in the prior formula whose inclusion restores equivalence. Around the same time, the IAV model was revived in practical computations. Potel et al.\ [Eur.\ Phys.\ J.\ A 53, 178 (2017)] benchmarked three independent IAV codes and applied them to deuteron-induced reactions on Ca isotopes with a dispersive optical model for the neutron-target propagator. That work demonstrated the predictive power of the IAV formalism for nonelastic breakup (NEB) cross sections and its applicability to exotic nuclei far from stability, establishing the NEB component as a key observable in inclusive breakup experiments.
